\section{Initial list gates and resampling estimates}
\label{sec:initial-list-gates-and-resampling-estimates}

The fresh cell sampler, used on typical pools with its required history
and bin conditions, makes every internal prior valid. We now require
that the early labels in other cells leave a substantial common-neighbor
list at each even role. After proving a failure estimate conditional on
typical pools, we will resample cells for finitely many rounds. The final
states must also retain a multiplicative comparison with fresh states;
this comparison will later be applied to products of list weights.

Let \(B_v^e\) be the external early neighbors of an even role \(v\).
Here external means noninternal, including both bulk and crossing edges.
Using the initial cell states, define the unnormalized row and its mass by
\[
 U_v^0(x)=\sigma_v(x)\prod_{b\in B_v^e}
                   \frac{I_x(y_b)}{D_{i(b)}(x)},\qquad
 M_v^0=\sum_x U_v^0(x).
\]
The internal prior includes hits on the initial dummy labels at any late
internal neighbors. For \(J\subseteq B_v^e\), let
\[
 L_v(J)=\{x\in\mathcal X_{i(v)}:
                  I_x(y_b)=1\text{ for every }b\in B_v^e\setminus J\}.
\]
This unweighted set imposes no internal hits. The event to be avoided is
\[
 S_v=\{M_v^0<1/2\}\ \cup\
 \left\{\max_{\substack{J\subseteq B_v^e\\|J|\le(\log n)^4}}
                       |L_v(J)|>\exp((\log n)^8)\right\}.
\]
The second clause keeps a small set of possible even labels after selected
early data are withheld. This is the support bound used in the late-exposure
comparison of Section~\ref{sec:transfer-for-a-late-prefix}.

\begin{lemma}[Independent list estimate with pins]\label{lem:independent-pinned-list}
Fix \(s\le s_0:=\lceil20R\rceil\) distinct members of \(B_v^e\) and
arbitrary labels for them. Also fix a valid probability prior \(\sigma_v\)
whose mass on the intersection of these pinned neighborhoods is at least
\(.9\cdot2^{-s}\). Draw the remaining external early labels independently
from their \(\pi_{i(b)}\)'s. Then \(\Pr(S_v)\le n^{-2R}\).
\end{lemma}

\Needspace{5\baselineskip}
\begin{proof}\leavevmode\par
\paragraph{Retained mass.}
Write \(i=i(v)\). Partition the external early incidences into the \(s\)
pins, the \(c'\) unpinned crossings, and the \(d\) unpinned bulk edges.
There is at most one late internal neighbor, so
\[
 s+c'+d=n-h_i-\#\{\text{external late neighbors}\}
             \le n-h_i-p(v)+1,
 \qquad d=n-O(\log n).
\]
Applying the pinned hit factors and their denominators first leaves mass
at least \(.9-o(1)\), since each denominator is \(.5+o(1)\) and \(s\)
is fixed. Normalize this remaining row, then expose the unpinned crossing
labels successively. At each step deep discrepancy
\eqref{eq:source-9} gives hit mass \(.5\pm O(b_*)\) except with
probability exponentially small in a power of \(n\). The intermediate
normalized laws remain small-width, and the cleaned crossing denominators
have the same estimate. Since \(c'b_*=o(1)\), the mass before the bulk
steps exceeds \(.8\) outside this exception. Its normalized law \(\tau\)
satisfies
\[
 \tau\le K2^s4^{c'}\sigma_v.
\]

The remaining bulk labels all have law \(\pi_i\). For the homogeneous
peeling estimate their parameter obeys
\[
 \Gamma:=\Delta_i\max_x\bigl(\tau(x)D_i(x)^{-d}\bigr)
          \le K e^{-200s_i}2^{-p(v)}.
\]
In cluster mode this follows by combining
\(N\|\sigma_v\|_\infty\le2^{h_i}e^{-500s_i}\),
\(\Delta_i\le e^{s_i}\), the degree-drift bound
\eqref{eq:source-12}, and \(c'\le\ell_i\le s_i/(1000u)\).
The preceding edge count and \(N\ge2^n\) give the factor
\(2^{-p(v)}\), with the possible extra factor 2 incorporated into \(K\).
In direct bias mode the positive bulk degree gain absorbs the patch-width,
conflict and crossing factors. In the bounded case those factors are fixed,
so the same inequality holds with a fixed \(K\) and \(s_i=0\).
Finally
\(p(v)\ge r/2\ge A_0\log n/2\); the chosen \(A_0\gg R\) therefore
makes \(\Gamma\le n^{-3R}\) for large \(n\).

Let \(Z\) be the bulk mass starting from \(\tau\). The centered-moment
identity, the moderate-interaction bound and homogeneous peeling give
\[
 \mathbb E|Z-1|^u\le n^{-3R}+O_u(\Gamma),\qquad
 \Pr(Z<3/4)\le4^u\mathbb E|Z-1|^u.
\]
Together with the crossing exception this is below the required
\(n^{-2R}\) bound with room to spare. On its complement,
\(M_v^0\ge .8(3/4)>1/2\).

\paragraph{Support size after omissions.}
For a fixed \(J\), ignore all internal hits, crossings and pins. The
remaining unpinned bulk hits are independent. The bound
\(n|D_i-.5|=O(\log n)\), the edge count above, and \(N\le n2^n\)
give
\[
 \mathbb E|L_v(J)|\le\exp(O((\log n)^5)).
\]
There are at most \(\exp(O((\log n)^5))\) deletion sets with
\(|J|\le(\log n)^4\). Markov at \(\exp((\log n)^8)\) and a union
over these sets give a probability much smaller than \(n^{-2R}\).
Adding it to the mass estimate proves the lemma.
\end{proof}

\subsection[Pool estimates]{Pool estimate for \(S_v\)}
\label{sec:pool-estimate-for-s-v}

At one star, fresh sampling uses the own cell of \(v\) for its internal
prior and at most one external early label from each other involved cell.
These cells are distinct: bulk neighbors lie in different nearby slices,
which cannot share a cell by the slice separation, and crossings visit
distinct other patches.

\begin{lemma}[Uniform pool estimate]\label{lem:uniform-pool-list-estimate}
Put \(T_s=\lceil\log^2 n\rceil\). Except on a pool event of probability
at most \(n^{-RT_s/2}\) per star, all involved individual pools are
typical and
\begin{equation}
 p_S(v\mid\text{pools})
 :=\Pr_{\rm fresh\ cells}(S_v\mid\text{pools})\le n^{-P}.          \label{eq:source-23}
\end{equation}
The same exception bound holds conditional on any one prescribed global
slot-to-bin pin.
\end{lemma}

\begin{proof}\leavevmode\par
\paragraph{Compatibility with a few prescribed slot labels.}
In addition to individual typicality, define the pool event
\(\mathcal A_v\) as follows. For every choice of at most \(s_0\)
distinct members of \(B_v^e\), one slot in each of their cells, and one
label in each selected bin permitted at that member, require
\[
 \Pr_{\text{fresh own-cell tape}}
 \left(\sigma_v\Bigl(\bigcap_{b\ {\rm chosen}}N_G(y_b)\Bigr)
                       <.9\,2^{-s}\ \middle|\ \text{own pool}\right)
       \le n^{-2R},
\]
where \(s\) is the number of chosen members and
\(N_G(y)=\{x:I_x(y)=1\}\). This is a predicate of the
pools: it bounds a conditional probability over a later own-cell draw.

Its failure, on individual typicality, has probability
\(\exp(-n^{\Omega(1)})\). To prove this, fix members, slot indices and
within-bin label indices, and compare the involved permutation slots to
iid slots, retaining a possible global pin. There are three cases. If the
pinned slot is one of the selected external slots, expose its label first.
The permission test \eqref{eq:source-21} and
Lemma~\ref{lem:fresh-internal-prior}, applied to the unforced own pool,
give an exponentially small averaged probability of a bad first degree.
The other selected slots are unforced. If the pin lies in the own cell,
all selected external slots are unforced, and only validity and the width
bound of the own prior are needed. If the pin lies in neither place, the
same argument uses entirely unforced selected data.

A prescribed within-bin index in a uniform bin has small width. Hence at
each unforced slot, \eqref{eq:source-9} puts its degree into the current
normalized first law in \(.5\pm2b_*\), except with probability
exponentially small in a power of \(n\). Successful restrictions preserve
small width through these at most \(s_0\) steps, and their intersection
mass is at least
\((.5-2b_*)^s\ge .9\,2^{-s}\). Thus the incompatibility probability,
averaged over the iid pools and the own tape with the needed typicality
indicators, is exponentially small. Markov at \(n^{-2R}\), followed by
a union over the \(n^{O(s_0(A_c+1))}\) choices of members, slots and
within-bin indices, proves the asserted bound for \(\mathcal A_v\).

\paragraph{Independent trials sharing their pools.}
Let \(\mathcal T_v\) require individual typicality and \(\mathcal A_v\).
Take \(T_s\) independent fresh-cell trials at these same pools. Then
\[
 \mathbb E_{\mathcal P}
       [\mathbf1_{\mathcal T_v}p_S(v\mid\mathcal P)^{T_s}]
 =\Pr(\mathcal T_v\text{ and every fresh trial has }S_v).
\]
Use the permutation-to-iid upper comparison, maintaining a global pin if
present. In each trial, \eqref{eq:source-22} bounds each external
singleton by a uniform choice of a slot index, with subsequent bin-label
weight
\[
 B_{i(b)}\pi_{i(b)}(y)
       \mathbf1_{\{D(y)=D_{\rm slot},\ {\rm permitted}\}}.
\]
Indeed averaging this weight over the \(L_{i(b)}^{\rm slot}\) indices
recovers the right side of \eqref{eq:source-22} without its error factor.
The accumulated factor over at most \(nT_s\) singletons is \(1+o(1)\).

There are now three remaining layers of randomness: selected slot indices,
the bin images of those slots, and the within-bin labels under the displayed
weights. Let \(j\) be the number of selected slot occurrences minus the
number of distinct unpinned slots visited. The fraction of index choices
with overlap rank \(j\) is at most
\[
 (2nT_s)^{3j}\bigl(\min_i L_i^{\rm slot}\bigr)^{-j}.
\]
One obtains this bound by specifying which choices repeat earlier indices
or hit the pinned index. A repeated unpinned slot of multiplicity \(q\)
contributes \(q-1\) to \(j\), and its \(q\) occurrences number at most
\(2(q-1)\). A pinned slot contributes its entire multiplicity to \(j\).
Thus at most \(2j\) occurrences use repeated or pinned slots. Call their
labels the trial pins.

Fix the images of those slots, their trial-pin labels, and the own pool.
Each such occurrence has total bin-label weight at most 11, because
\(\pi_i\le11/M_i\) and a bin contains \(d_i\) labels; thus all trial
pins cost at most \(11^{2j}\). The unique unforced slot images are still
independent. Integrating each restores \(\pi_{i(b)}\), after discarding
its permission indicator for an upper bound. For a trial with at most
\(s_0\) pins, retain from \(\mathcal A_v\) only its own-tape
compatibility assertion. This assertion does not read any of the unique
unforced slots. Its incompatible own-tape probability is at most
\(n^{-2R}\); on compatibility, Lemma~\ref{lem:independent-pinned-list}
gives another \(n^{-2R}\). Hence the trial's bad probability is at most
\(n^{-R}\). The fresh own tapes and the remaining unforced slots are
independent across trials, so these bounds multiply.

At most \(2j/s_0\) trials have too many pins and require the trivial
bound 1. Summing overlap ranks therefore gives
\[
 \begin{aligned}
 \mathbb E[\mathbf1_{\mathcal T_v}p_S^{T_s}]
 &\le O(1)n^{-RT_s}
    \sum_{j\ge0}
    \left[
      \frac{(2nT_s)^3\,11^2\,n^{2R/s_0}}
           {\min_i L_i^{\rm slot}}
    \right]^j\\
 &\le O(1)n^{-RT_s}.
 \end{aligned}
\]
Here \(2R/s_0\le .1\) and
\(L_i^{\rm slot}=n^{A_c+o(1)}\), with \(A_c=200\), so the bracket
tends to zero. This proves the moment bound with its typicality indicator.
Finally
\[
 \Pr(p_S>n^{-P},\mathcal T_v)
 \le n^{PT_s}\mathbb E[\mathbf1_{\mathcal T_v}p_S^{T_s}]
 \le O(1)n^{-(R-P)T_s}.
\]
Since \(R=P^2\), this and the exponentially smaller individual and
auxiliary pool exceptions are at most \(n^{-RT_s/2}\) for large \(n\).
\end{proof}

\subsection{Palettes}
\label{sec:palettes}

\begin{lemma}[Palette retention and pair estimates]\label{lem:low-mode-palettes}
There is a fixed palette assignment with the mass retention and separation
properties below, uniform over every valid local history. It satisfies
\eqref{eq:source-24} and the accompanying second-moment bound.
\end{lemma}

\begin{proof}\leavevmode\par
\paragraph{Separating internal words.}
Choose a binary linear map on the internal words of patch \(i\) whose
kernel contains no nonzero word of weight at most \(500\rho h_i\).
Random linear checks and the ball-volume bound give such a map with image
size \(\chi_i\le e^{s_i}\). Set \(\chi_i=1\) when there are no internal
bits. Its colors separate distinct words at the indicated distance.
Among even roles the colors have equal sizes, since a free outer bit
supplies either parity for each internal word.

\paragraph{Retaining list mass.}
Independently color the labels of \(X_i\) uniformly by these \(\chi_i\)
colors. We can fix a realization with each palette of size at most
\(2M_i/\chi_i\), such that for every possible valid local history with
\(M_v^0\ge1/2\), its assigned palette satisfies
\[
 \sum_{x\ {\rm in\ the\ palette\ of}\ v}U_v^0(x)
       \ge 1/(4\chi_i).
\]
At the same time, every palette has \(\sigma_v\)-mass at most
\(2/\chi_i\) for every possible valid internal prior.

For this simultaneous assertion the deterministic atom estimate is
\[
 \max_x U_v^0(x)\le K C_n^{-1}e^{-200s_i}2^{-p(v)}.
\]
It follows from the cap and degree calculation in the pinned-list proof.
The logarithm of the number of histories is at most \(n^{2+o(1)}\):
Lemma~\ref{lem:finite-low-mode-data} counts the own slice data, and the
star labels add at most \(n\log N=O(n^2)\). Weighted-sum concentration
for the random label colors has exponent at least a fixed multiple of
\[
 C_n e^{199s_i}2^{p(v)},
\]
using \(\chi_i\le e^{s_i}\) and \(M_v^0\ge1/2\). The bound
\(p(v)\ge A_0\log n/2\) makes this dominate the logarithmic history
count. The much smaller atoms of \(\sigma_v\) give the simultaneous
upper palette-mass bound, and ordinary size concentration gives the
palette sizes. Thus one coloring has all the asserted properties.
After restriction to the assigned palette and normalization, the starting
row has cap at most
\(4K C_n^{-1}e^{-199s_i}2^{-p(v)}\).

\paragraph{Pair estimates.}
For all these valid internal priors and all \(x\in\mathcal X_i\),
\begin{equation}
 \sum_{\substack{z\ {\rm in\ assigned\ palette}\\
                  |K_{\pi_i}(x,z)|\le\xi}}
  \sigma_v(z)\prod_{b\in B_v^e}
    \frac{\int I_xI_z\,d\pi_{i(b)}}{D_{i(b)}(x)D_{i(b)}(z)}
       \le K/\chi_i .                                          \label{eq:source-24}
\end{equation}
To verify this, write the ratio for a law \(\pi_j\) as
\[
 R_j(x,z)=1+
 \frac{K_{\pi_j}(x,z)-m_{\pi_j}(x)m_{\pi_j}(z)}
      {(1+m_{\pi_j}(x))(1+m_{\pi_j}(z))}.
\]
First consider pairs typical in every relevant relaxed row-tail test
\eqref{eq:source-11}. Each bulk ratio is
\(1+O(n^{-1-.02}+\log^2n/n^2)\); each crossing ratio is
\(1+o(1/\log n)\). There are at most \(n\) bulk factors and
\(\ell_i=o(\log n)\) crossings, so their product is bounded. The
palette's \(\sigma_v\)-mass is at most \(2/\chi_i\), proving the
claim on this part of the sum.

For bulk exceptions within the nonconflict cutoff, the weighted row tail
in relaxed \eqref{eq:source-11} controls the product of bulk ratios.
Its \(e^{-n^{.19}}\) bound also absorbs the crude
\(\exp(O(\ell_i))\) crossing factor. For crossing exceptions with
typical bulk, the unweighted row tail suffices. In each case replacing
uniform \(z\in X_i\) by \(\sigma_v\) costs only
\(\exp(O(\log n))\), by the low-mode cap. The resulting exception
contribution is negligible compared with \(1/\chi_i\), proving
\eqref{eq:source-24}.

We also record the precise uniform integral used later. Let
\(\mathcal A(x,z)\) mean that both labels belong to \(\mathcal X_i\)
and are nonconflicting. Then, for \(0\le d\le n\),
\[
 \mathbb E_{x,z\ {\rm iid\ uniform\ on}\ X_i}
 \left[\mathbf1_{\mathcal A(x,z)}
       \left(4\int I_xI_z\,d\pi_i\right)^{2d}\right]
       \le\exp(O(\log n)).
\]
Indeed the degree terms contribute \(\exp(O(\log n))\); the typical
pair term is bounded, and the same weighted tail controls the exceptions
inside the cutoff. This is an unconditioned uniform integral with an
indicator, rather than a law conditioned on the envelope or a palette.
\end{proof}

\subsection{Parallel initial resampling}
\label{sec:parallel-initial-resampling}

We use the independent sample tapes and backward execution witnesses of
the Moser--Tardos resampling method~\cite{moser-tardos2010}.
Haeupler, Saha and Srinivasan~\cite[Theorem~2.2]{haeupler-saha-srinivasan2011}
established multiplicative comparisons for fixed events determined by the
current assignment during a resampling execution; Harris and
Srinivasan~\cite{harris-srinivasan2017} developed this distributional
analysis further. The parallel-round theorem of Moser and Tardos uses a
maximal independent set in each round. The fixed-priority rule below need
not choose a maximal independent set, so we prove its finite-round and
final-state estimates directly.

Fix the pools, and give each cell its independent tape of fresh states.
The scope of \(S_v\) is the set of cells it reads. Form the scope graph
with one vertex per event \(S_v\), joining events whose scopes share a
cell, and fix an ordering of its vertices. At each of \(T_s\) rounds,
execute every currently true event having no earlier-numbered true
neighbor. An execution advances the tape index of every cell in its scope.
The executing scopes are disjoint, although some true events may have no
executing neighbor. Scope sizes, cell incidences and graph degrees are at
most \(n^{A_c+4}\).

A restricted simulation uses a deterministic subset of events, its induced
scope graph, and the same priorities and tapes. Truth outside that subset
is not consulted. During one round, update and priority influences travel
at most two graph steps. Hence the final state of a target cell, or the
final truth of a target \(S_v\), agrees with the global simulation whenever
the restricted set contains the graph ball of radius \(2T_s+3\) about
that cell's incidence events or about \(v\), respectively. The same
holds for any larger deterministic set.

\begin{proposition}[Finite resampling and multiplicative comparison]\label{prop:initial-resampling-comparison}
Fix, before drawing tapes, a deterministic restricted event set, a defining
event \(S_v\) in that set, and a deterministic set of at most
\(n^{10(A_c+10)}\) target cells. Suppose the pools are individually
typical and satisfy \eqref{eq:source-23} for every event tested in the
restricted simulation. All conclusions below concern that restricted
process. The defining event's final failure probability is at most
\(n^{-PT_s/2}\). The probability of more than \(T_s\) executions in the
backward closure of the target cells is also at most \(n^{-PT_s/2}\).
For every nonnegative test of those target states, the final-state
expectation is at most \(1+o(1)\) times its independent fresh-cell
expectation at the same fixed pools.
\end{proposition}

\Needspace{5\baselineskip}
\begin{proof}\leavevmode\par
All probabilities in this proof condition on the fixed pools. Every
fresh truth test has probability at most \(n^{-P}\).

\paragraph{A witness for final failure.}
If the defining event \(S_v\) is true at the finish, consider its
component in the restricted scope graph induced by all sites ever true,
at round starts or at the finish. This component has a true site at every
round start. Indeed, if it had none at some round, no outside execution
could change any of its truth values: an executing outside site would be
true and adjacent to it, and hence belong to the same component. It would
then stay entirely false, contradicting final truth. The least-numbered
true site in the component therefore executes each round.

List all execution occurrences in this component, say \(M\ge T_s\).
Add a maximal collection of \(J\) further ever-true sites nonadjacent to
any execution site and pairwise nonadjacent. The added sites have untouched
inputs. The distance-one neighborhoods of the listed sites cover the
component; consequently the listed sites are connected by graph steps of
length at most three and anchored near \(v\). Rooted tree traversals,
with occurrence multiplicities, types and execution rounds specified,
bound the number of such lists by
\[
 \bigl(n^{O(A_c+4)}T_s^{O(1)}\bigr)^{M+J}.
\]

For each possible list, determine the tape entries used by its truth
tests from the listed prior executions. Every execution touching one of
these sites belongs to the same component and has been listed. The event
that this list is the actual witness is contained in the event that all
these prescribed truth tests pass. These necessary tests use independent
entries: an execution consumes each entry it tests, while the extra sites
read untouched, mutually disjoint scopes. Their joint probability is at
most \(n^{-P(M+J)}\), before any consistency conditions on the list are
imposed. The enumeration and this probability form a geometric sum with
per-occurrence factor
\(n^{-P+C(A_c+4)+o(1)}\). The sufficiently large choice of \(P\),
and \(M\ge T_s\), bound the sum by \(n^{-PT_s/2}\).

\paragraph{Backward executions and terminal entries.}
For the fixed target cells, list every execution touching them, then
recursively include every earlier execution whose scope meets a listed
later one. This is their backward closure. Its occurrence lists are
counted in the same way, joining all target incidences to one auxiliary
root. The polynomial number of targets is included in the constant in
\(n^{O(A_c+4)}\). A list of size \(m\) has at most
\(\bigl(n^{O(A_c+4)}T_s^{O(1)}\bigr)^m\) possibilities, and its
necessary truth tests have probability at most \(n^{-Pm}\). Summing
\(m>T_s\) gives the second assertion.

For a fixed occurrence list, the final entry of each target cell is the
entry immediately after its last listed touch, or its initial entry if
there was no touch. These indices are prescribed by the list. None is
used in a necessary truth test: every execution touching a target has
been included, and all its consumed entries precede the terminal one.
The terminal entries are thus independent fresh states, independent also
of the necessary truth tests. For any nonnegative target test \(\Psi\),
consistency of the list may be discarded, leaving
\[
 \begin{aligned}
 \mathbb E[\Psi(\text{final targets})]
 &\le\mathbb E[\Psi(\text{fresh targets})]
    \left(1+\sum_{m\ge1}
           [n^{-P+C(A_c+4)+o(1)}]^m\right)\\
 &=(1+o(1))\mathbb E[\Psi(\text{fresh targets})].
 \end{aligned}
\]
The first term accounts for an empty closure. This factorization uses
only the pool requirements in the specified deterministic horizons and
proves the multiplicative comparison.
\end{proof}

\begin{remark}[Scope of the comparison]\label{rem:resampling-test-scope}
The target cells and simulation horizons in
Proposition~\ref{prop:initial-resampling-comparison} are fixed before
their tapes are drawn. No buffer is needed for the restricted-process
estimates; the preceding radius condition is used only when identifying
these states with those of the global run. The comparison applies directly to nonnegative
list products, regardless of their size. An additive rare-history bound
would not give this conclusion: the proof instead separates the consumed
truth-test entries from the fresh terminal entries before discarding
occurrence-pattern consistency.
\end{remark}
